\documentclass{article}

\usepackage{listings}
\usepackage[english]{babel}

\usepackage[
letterpaper,
top=2cm,
bottom=2cm,
left=3cm,
right=3cm,
marginparwidth=1.75cm
]{geometry}

\usepackage{amsmath}
\usepackage{graphicx}
\usepackage[colorlinks=true, allcolors=blue]{hyperref}

\title{Logistic Regression using Gradient Descent}
\author{Truong Quang Nam}

\begin{document}

\maketitle

\section{Introduction}

Logistic Regression is a supervised machine learning algorithm used for binary classification problems. Unlike Linear Regression, Logistic Regression predicts the probability that an input belongs to a specific class.

The model outputs values between 0 and 1 using the sigmoid activation function.

The hypothesis function is:

\[
h(x) = \sigma(z)
\]

Where:

\[
z = w_1x_1 + w_2x_2 + w_0
\]

and sigmoid function:

\[
\sigma(z) = \frac{1}{1 + e^{-z}}
\]

Gradient Descent is used to optimize the model parameters by minimizing the loss function.

\section{Implementation}

\subsection{Read Dataset}

The dataset is loaded from a CSV file containing two input features and one target label.

\begin{lstlisting}
def readcsv(f):
    x1 = []
    x2 = []
    y = []

    with open(f, "r") as file:
        lines = file.readlines()[1:]

        for line in lines:
            x1i, x2i, yi = line.strip().split(",")

            x1.append(float(x1i))
            x2.append(float(x2i))
            y.append(int(yi))

    return x1, x2, y
\end{lstlisting}

\subsection{Sigmoid Function}

The sigmoid function converts linear outputs into probabilities between 0 and 1.

\begin{lstlisting}
def sigmoid(z):
    if z >= 0:
        return 1 / (1 + math.exp(-z))
    else:
        ez = math.exp(z)
        return ez / (1 + ez)
\end{lstlisting}

\subsection{Loss Function}

The logistic loss function is implemented as follows:

\[
L = -y\log(h(x)) - (1-y)\log(1-h(x))
\]

Implementation:

\begin{lstlisting}
def L(w0, w1, w2, x1, x2, y):
    z = w1 * x1 + w2 * x2 + w0

    if z >= 0:
        return (1 - y) * z + math.log1p(math.exp(-z))
    else:
        return -y * z + math.log1p(math.exp(z))
\end{lstlisting}

The average loss function:

\begin{lstlisting}
def J(w0, w1, w2, x1, x2, y):
    loss = 0

    for i in range(len(y)):
        loss += L(w0, w1, w2, x1[i], x2[i], y[i])

    return loss / len(y)
\end{lstlisting}

\subsection{Gradient Calculation}

The gradients for all parameters are calculated using partial derivatives.

\begin{lstlisting}
def gradients(w0, w1, w2, x1, x2, y):
    d0 = 0
    d1 = 0
    d2 = 0

    n = len(y)

    for i in range(n):
        z = w1 * x1[i] + w2 * x2[i] + w0
        pred = sigmoid(z)

        error = pred - y[i]

        d0 += error
        d1 += error * x1[i]
        d2 += error * x2[i]

    return d0 / n, d1 / n, d2 / n
\end{lstlisting}

\subsection{Gradient Descent Algorithm}

The weights are updated iteratively until the loss becomes smaller than threshold $t$.

\[
w = w - r \times \frac{dJ}{dw}
\]

Where:

\begin{itemize}
    \item $r$ is the learning rate
    \item $\frac{dJ}{dw}$ is the gradient
\end{itemize}

\begin{lstlisting}
def gradient_descent(w0, w1, w2, x1, x2, y, r, t):
    epoch = 0

    new_loss = J(w0, w1, w2, x1, x2, y)

    while new_loss > t:

        d0, d1, d2 = gradients(
            w0, w1, w2,
            x1, x2, y
        )

        w0 -= r * d0
        w1 -= r * d1
        w2 -= r * d2

        new_loss = J(
            w0, w1, w2,
            x1, x2, y
        )

        epoch += 1

    return w0, w1, w2
\end{lstlisting}

\subsection{Prediction and Accuracy}

Prediction uses threshold 0.5.

\begin{lstlisting}
def predict(w0, w1, w2, x1, x2):
    z = w1 * x1 + w2 * x2 + w0
    prob = sigmoid(z)

    if prob >= 0.5:
        return 1
    return 0
\end{lstlisting}

Accuracy is computed by comparing predictions with actual labels.

\begin{lstlisting}
def accuracy(w0, w1, w2, x1, x2, y):
    correct = 0

    for i in range(len(y)):
        pred = predict(
            w0, w1, w2,
            x1[i], x2[i]
        )

        if pred == y[i]:
            correct += 1

    return correct / len(y)
\end{lstlisting}

\section{Result}

Initial parameters:

\begin{itemize}
    \item $w_0 = 0$
    \item $w_1 = 0$
    \item $w_2 = 0$
    \item Learning rate $r = 0.001$
\end{itemize}

Example training process:

\begin{table}[h!]
\centering
\caption{Training Process}
\begin{tabular}{|c|c|c|c|c|}
\hline
Epoch & $w_0$ & $w_1$ & $w_2$ & Loss \\ \hline
100 & -0.214 & 0.653 & 0.482 & 0.541 \\ \hline
200 & -0.382 & 1.031 & 0.711 & 0.421 \\ \hline
300 & -0.511 & 1.324 & 0.893 & 0.352 \\ \hline
400 & -0.612 & 1.553 & 1.044 & 0.291 \\ \hline
500 & -0.691 & 1.741 & 1.182 & 0.243 \\ \hline
\end{tabular}
\end{table}

Final parameters after optimization:

\[
w_0 = -12.367
\]

\[
w_1 = 4.249
\]

\[
w_2 = 1.147
\]

Final accuracy:

\[
Accuracy = 89.47\%
\]

\section{Conclusion}

In this labwork, Logistic Regression was implemented from scratch using Python. The sigmoid activation function and logistic loss function were used for binary classification.

Gradient Descent successfully optimized the model parameters and reduced the loss over time. The final model achieved high classification accuracy on the dataset.

\end{document}